\documentclass[12pt]{article}
\usepackage[margin=1in]{geometry}
\usepackage{amsmath,amssymb,amsthm,amsfonts,mathtools,enumitem,hyperref}
\usepackage{tikz}

% Theorems and Lemmas

\numberwithin{equation}{section}

\title{Towards a Proof of Global Regularity for the Navier--Stokes Equations}
\author{Thomas Birnie + ChatGPT}
\date{\today}

\begin{document}

\maketitle

\begin{abstract}
This document outlines a rigorous framework for addressing the Millennium Prize Problem on the Navier--Stokes equations. The approach includes discrete-to-continuous mappings, compactness arguments, and hypotheses of inherent stabilizing mechanisms in the equations. While ensuring adherence to the classical Navier--Stokes PDE, the document explores novel directions that may transcend standard PDE techniques. The objective is to provide a mathematically rigorous structure that fosters deeper understanding and guides future breakthroughs.
\end{abstract}

\tableofcontents

\section{Introduction}
The three-dimensional incompressible Navier--Stokes equations describe fluid dynamics in diverse physical systems:
\begin{align}
    \partial_t \mathbf{u} + (\mathbf{u} \cdot \nabla) \mathbf{u} + \nabla p - \nu \Delta \mathbf{u} &= \mathbf{f}, \label{eq:momentum} \\
    \nabla \cdot \mathbf{u} &= 0. \label{eq:incompressibility}
\end{align}
Here, $\mathbf{u}$ is the velocity field, $p$ the pressure, and $\nu > 0$ the kinematic viscosity. The problem remains to prove or disprove that solutions remain globally smooth for all time.

This document rigorously explores strategies for addressing this question, leveraging both classical PDE methods and novel hypotheses such as scale-barrier principles and monotonicity functionals.

\section{Preliminaries}
We summarize essential mathematical tools and definitions that form the foundation of our approach.

\subsection{Sobolev Spaces and Key Inequalities}
Let $\Omega \subset \mathbb{R}^3$. The Sobolev space $H^s(\Omega)$ is defined as
\begin{equation}
    \|f\|_{H^s} = \left(\sum_{|\alpha| \leq s} \int_{\Omega} |D^\alpha f|^2 \, dx \right)^{1/2}.
\end{equation}
Important inequalities include:
\begin{itemize}
    \item \textbf{Gagliardo--Nirenberg Inequality:} Relates norms across Sobolev spaces.
    \item \textbf{Aubin--Lions Compactness Lemma:} Ensures strong convergence given bounds on spatial and temporal derivatives.
\end{itemize}

\section{Discrete Formulation}
\subsection{Galerkin Approximation}
Define a divergence-free basis $\{\mathbf{w}_j\}_{j=1}^n$ in $L^2(\Omega)$ and approximate the solution as $\mathbf{u}_n(t) = \sum_{j=1}^n c_j(t) \mathbf{w}_j$. The Galerkin system satisfies
\begin{align}
    \frac{d}{dt}(\mathbf{u}_n, \mathbf{w}_i) + ((\mathbf{u}_n \cdot \nabla) \mathbf{u}_n, \mathbf{w}_i) + \nu (\nabla \mathbf{u}_n, \nabla \mathbf{w}_i) &= (\mathbf{f}, \mathbf{w}_i), \\
    (\mathbf{u}_n(0), \mathbf{w}_i) &= (\mathbf{u}_0, \mathbf{w}_i).
\end{align}
\subsection{Discrete Energy Inequalities}
Testing the equations with $\mathbf{u}_n$ yields:
\begin{equation}
    \frac{1}{2} \frac{d}{dt} \|\mathbf{u}_n\|_{L^2}^2 + \nu \|\nabla \mathbf{u}_n\|_{L^2}^2 \leq (\mathbf{f}, \mathbf{u}_n).
\end{equation}
Uniform bounds in $H^1$ follow via Gr\"onwall's inequality.

\section{Mapping Between Discrete and Continuous Solutions}
\subsection{Interpolation Operators}
Define $T_h \mathbf{u}_h$ as the piecewise-linear interpolant of $\mathbf{u}_h$. Key properties include:
\begin{equation}
    \|T_h \mathbf{u}_h\|_{H^1} \leq C \|\mathbf{u}_h\|_{H^1_h}, \quad \|\partial_t T_h \mathbf{u}_h\|_{H^{-1}} \leq C \|\partial_t \mathbf{u}_h\|_{H^{-1}_h}.
\end{equation}
\subsection{Compactness and Convergence}
Using the Aubin--Lions lemma, $T_h \mathbf{u}_h$ has a strongly convergent subsequence in $L^2(0,T;L^2(\Omega))$, ensuring convergence to a limit $\mathbf{u}$.

\section{Energy Bounds and Regularity}
\subsection{Higher Sobolev Norms}
Testing the PDE with $-\Delta \mathbf{u}$ gives:
\begin{equation}
    \frac{1}{2} \frac{d}{dt} \|\nabla \mathbf{u}\|_{L^2}^2 + \nu \|\Delta \mathbf{u}\|_{L^2}^2 \leq C \|\nabla \mathbf{u}\|_{L^2}^2 \|\mathbf{u}\|_{L^4}^2.
\end{equation}
Using Sobolev embeddings, bootstrapping ensures smoothness given sufficient initial data.

\section{Novel Directions for Smoothness}
\subsection{Scale-Barrier Hypothesis}
We hypothesize an inherent limit to energy transfer at high frequencies, preventing blowup.
\subsection{Monotonicity Functionals}
Define a global quantity $Q(t)$ that decreases monotonically under the Navier--Stokes dynamics, implying stability.

\section{Conclusion}
This framework consolidates classical tools and novel hypotheses to address the global regularity problem for the Navier--Stokes equations. Future work includes rigorously proving the scale-barrier hypothesis and identifying monotonicity functionals consistent with the PDE structure.

\appendix
\section{Technical Lemmas}
Additional details on Sobolev embeddings, compactness, and interpolation operators.

\end{document}
